\documentclass[12pt]{article}

% ============================================================
%  PAQUETES
% ============================================================
\usepackage[spanish]{babel}

% ------------------------------------------------------------
%  TIPOGRAFÍA: Arial en el texto, monoespaciada en el código
%  IMPORTANTE: este archivo se compila con XeLaTeX o LuaLaTeX,
%  no con pdfLaTeX. En Overleaf: Menu > Compiler > XeLaTeX.
%  Si Arial no está instalado se usa Liberation Sans, que tiene
%  exactamente las mismas métricas y no altera la paginación.
% ------------------------------------------------------------
\usepackage{fontspec}
\IfFontExistsTF{Arial}{%
  \setmainfont{Arial}%
  \setsansfont{Arial}%
}{%
  \setmainfont{Liberation Sans}%
  \setsansfont{Liberation Sans}%
}
\IfFontExistsTF{Consolas}{%
  \setmonofont{Consolas}[Scale=MatchLowercase]%
}{%
  \setmonofont{Liberation Mono}[Scale=MatchLowercase]%
}

\usepackage[a4paper, margin=2.5cm]{geometry}
\usepackage{graphicx}          % Imágenes
\usepackage{float}             % Anclaje de figuras con [H]
\usepackage{enumitem}          % Control de listas
\usepackage{needspace}         % Evita separar un bloque de su captura
\usepackage{tocloft}           % Ancho de la numeración en el índice
\setlength{\cftsecnumwidth}{2em}
\setlength{\cftsubsecnumwidth}{2.8em}
\usepackage{xcolor}            % Colores
\usepackage{listings}          % Bloques de código
\usepackage{fancyhdr}          % Encabezado y pie de página
\usepackage{titlesec}          % Estilo de títulos de sección
\usepackage{microtype}         % Mejor justificación / tipografía
\usepackage{setspace}          % Interlineado
\onehalfspacing                % Interlineado 1.5 en todo el documento
\usepackage[font=small, labelfont={bf,color=purple-uach}]{caption}
\usepackage[hidelinks]{hyperref} % Enlaces e índice clicable
\usepackage{xurl}
\urlstyle{same}

\captionsetup{font={small,stretch=1.0}}    % pies de figura a renglón sencillo

\graphicspath{{img/}{./}}

% ============================================================
%  COLORES
% ============================================================
\definecolor{purple-uach}{HTML}{371e49}

% ============================================================
%  DATOS DE LA TAREA
% ============================================================
\newcommand{\alumnoUno}{Osvaldo Hernández Juárez}
\newcommand{\matriculaUno}{385493}
\newcommand{\alumnoDos}{Diego Flores Verdad Grijalva}
\newcommand{\matriculaDos}{385660}
\newcommand{\alumnoTres}{Luis Felipe Domínguez Chávez}
\newcommand{\matriculaTres}{385500}
\newcommand{\materia}{Bases de Datos Avanzadas}
\newcommand{\docente}{José Saúl De Lira Miramontes}
\newcommand{\grupoCarrera}{Grupo 5k2}
\newcommand{\actividad}{Proyecto Segundo Parcial: Bases de Datos Distribuidas}
% Versión de la portada, con el salto de renglón donde se ve mejor.
\newcommand{\actividadPortada}{Proyecto Segundo Parcial:\\ Bases de Datos Distribuidas}
\newcommand{\fechaEntrega}{5 de noviembre de 2026} % ajustar el día de la entrega

\hypersetup{
    pdftitle={\actividad},
    pdfauthor={\alumnoUno, \alumnoDos, \alumnoTres},
    pdfsubject={\materia}
}

% ============================================================
%  ESTILO DE TÍTULOS DE SECCIÓN
% ============================================================
\titleformat{\section}
  {\large\bfseries\color{purple-uach}}{\thesection}{1em}{}
\titleformat{\subsection}
  {\normalsize\bfseries\color{purple-uach}}{\thesubsection}{1em}{}
\titlespacing*{\section}{0pt}{2.4ex plus 1ex minus .2ex}{1.2ex plus .2ex}

% ============================================================
%  PIE DE PÁGINA
% ============================================================
\pagestyle{fancy}
\fancyhf{}
\renewcommand{\headrulewidth}{0pt}
\renewcommand{\footrulewidth}{0pt}
\fancyfoot[R]{\thepage}

% ============================================================
%  ESTILO DE BLOQUES DE CÓDIGO
% ============================================================
\lstdefinelanguage{mongosh}{
    morekeywords={db,sh,rs,const,let,new,find,findOne,insertOne,insertMany,updateOne,updateMany,deleteOne,deleteMany,countDocuments,aggregate,runCommand,getSiblingDB,createCollection,drop,print,printjson},
    morestring=[b]",
    morestring=[b]',
    morecomment=[l]{//},
    sensitive=true
}

\lstdefinestyle{codeStyle}{
    backgroundcolor=\color{gray!8},
    commentstyle=\color{purple-uach!70},
    keywordstyle=\color{purple-uach}\bfseries,
    stringstyle=\color{teal},
    basicstyle=\setstretch{1}\ttfamily\footnotesize,
    breaklines=true,
    breakatwhitespace=false,
    frame=single,
    rulecolor=\color{gray!40},
    tabsize=2,
    captionpos=b,
    aboveskip=1.2ex,
    belowskip=1.2ex
}
\lstset{style=codeStyle, language=mongosh}
\lstdefinestyle{pyStyle}{style=codeStyle, language=Python, showstringspaces=false, morekeywords={self,None,True,False}}
\lstdefinestyle{plainStyle}{style=codeStyle, language={}, showstringspaces=false}


% ============================================================
%  CAPTURAS Y PENDIENTES
% ============================================================
% \captura{archivo}{pie de figura}{ancho}: si la imagen todavía no está en
% img/, se muestra un recuadro con el nombre del archivo que falta.
\newcommand{\captura}[3]{%
  \begin{figure}[H]
    \centering
    \IfFileExists{img/#1}{%
      \includegraphics[width=#3\linewidth]{#1}%
    }{%
      \fbox{\parbox[c][2.5cm][c]{0.8\linewidth}{\centering\small\color{gray}Falta la captura \texttt{img/#1}}}%
    }
    \caption{#2}
  \end{figure}}
% \pendiente{texto}: aviso visible de algo que falta; quitarlo antes de entregar.
\newcommand{\pendiente}[1]{%
  \par\noindent\fcolorbox{purple-uach}{purple-uach!6}{\parbox{\dimexpr\linewidth-2\fboxsep-2\fboxrule}{\small\textbf{Pendiente:} #1}}\par}

\lstdefinestyle{jsStyle}{style=codeStyle, language=mongosh, showstringspaces=false, morecomment=[s]{/*}{*/}, morestring=[b]`}

\emergencystretch=1.5em

% ============================================================
\begin{document}

% ---------------- PORTADA ----------------
\begin{titlepage}
\singlespacing
\centering
\noindent
\begin{minipage}[c]{0.5\linewidth}
    \raggedright
    \includegraphics[height=4.5cm]{logo-uach.png}
\end{minipage}%
\begin{minipage}[c]{0.5\linewidth}
    \raggedleft
    \includegraphics[height=4.5cm]{logo-fing.png}
\end{minipage}\\[1.2cm]
{\Large \textbf{Universidad Autónoma de Chihuahua}}\\[0.2cm]
{\large Facultad de Ingeniería}\\[0.2cm]
{\normalsize \grupoCarrera}\\[2cm]

{\color{purple-uach}\rule{\linewidth}{0.6mm}}\\[0.5cm]
{\huge \bfseries \linespread{1.25}\selectfont \actividadPortada\par}\vspace{0.5cm}
{\color{purple-uach}\rule{\linewidth}{0.6mm}}\\[1.5cm]

{\Large \materia}\\[0.3cm]
{\large Docente: \docente}

\vfill

\begin{flushright}
    \large
    \textbf{\alumnoUno}\\
    Matrícula: \matriculaUno\\[0.4cm]
    \textbf{\alumnoDos}\\
    Matrícula: \matriculaDos\\[0.4cm]
    \textbf{\alumnoTres}\\
    Matrícula: \matriculaTres
\end{flushright}

\vspace{1cm}
{\large \fechaEntrega}

\end{titlepage}


% ---------------- ÍNDICE ----------------
\tableofcontents
\newpage

% ============================================================
%  CAPTURAS: cada quien numera las suyas en su rango
%    01-09  Luis Felipe (cluster, datos y shard key)
%    10-19  Diego (replica sets y balanceo)
%    20-39  Osvaldo (CRUD con monitoreo)
%  Las imágenes van en reporte/img/. Si una falta, el PDF muestra un
%  recuadro con el nombre del archivo.
% ============================================================

% ============================================================
\section{Introducción}

En la actividad anterior se configuró un cluster de MongoDB con sharding sobre Docker y una colección pequeña de usuarios. En este proyecto el mismo esquema se aplicó a una base de datos real, sample\_mflix, cuya colección movies contiene 21\,349 películas. Cada shard se configuró como un replica set de tres miembros, se agregó a cada película un identificador numérico, id\_movie, que se usó como shard key de tipo hashed, y sobre la colección distribuida se realizaron operaciones CRUD, monitoreando después de cada una cómo quedan repartidos los documentos entre los shards. También se muestra cómo el balanceador reparte la carga cuando se agrega un shard nuevo a un cluster que ya tiene datos.

Los archivos del proyecto (Docker Compose, scripts de configuración, monitoreo, balanceo y operaciones CRUD) se encuentran en el repositorio \url{https://github.com/LF-Dguez/bases_de_datos_distribuidas_MongoDB}.

% ============================================================
\section{Objetivo}

Implementar soluciones para Bases de Datos Distribuidas.

% ============================================================
\section{Entorno de trabajo}

El proyecto se realizó en Windows 11 con Docker Desktop y Docker Compose v2. Los once nodos del cluster usan la imagen oficial \texttt{mongo:8.0}. Los scripts se escribieron en JavaScript y se ejecutaron con el \texttt{mongosh} que incluye la imagen, dentro del contenedor del router; la importación de los datos se hizo con \texttt{mongoimport}, que también viene en la imagen.

% ============================================================
\section{Configurar un Cluster en Mongo sobre Docker con los siguientes nodos: configsvr, shard1, shard2, shard3, Replica set, mongos router}

\subsection{Arquitectura}

El cluster está formado por once contenedores:

\begin{itemize}[itemsep=2pt, topsep=4pt]
    \item \texttt{configsvr}: config server. Guarda los metadatos del cluster, como la lista de shards y los rangos de datos (chunks) de cada uno. Es un replica set de un miembro.
    \item \texttt{shard1}, \texttt{shard2} y \texttt{shard3}: cada shard es un replica set de tres miembros (por ejemplo, \texttt{shard1-1}, \texttt{shard1-2} y \texttt{shard1-3}). Uno de ellos es el PRIMARY, que recibe las escrituras, y los otros dos son SECONDARY y mantienen una copia de los mismos datos. Así cada parte de la colección está guardada tres veces y el shard sigue funcionando si falla uno de sus miembros.
    \item \texttt{mongos}: router. Las aplicaciones se conectan a él y él envía cada operación al shard que corresponde. No guarda datos.
\end{itemize}

\subsection{Docker Compose}

\lstinputlisting[style=plainStyle]{../docker-compose.yaml}

Los puntos principales de la configuración son:

\begin{itemize}[itemsep=2pt, topsep=4pt]
    \item El config server se inicia con \texttt{--configsvr} y los miembros de los shards con \texttt{--shardsvr}. Con \texttt{--replSet} cada uno recibe el nombre de su replica set; los tres miembros de un shard comparten el mismo nombre.
    \item Los nueve miembros de los shards comparten la misma configuración base, definida una sola vez con la extensión \texttt{x-shard} y reutilizada con \texttt{\textless\textless: *shard}.
    \item Todos escuchan en el puerto 27017 dentro de su contenedor y aceptan conexiones de los demás con \texttt{--bind\_ip\_all}. Cada \texttt{mongod} limita su caché a 0.25 GB con \texttt{--wiredTigerCacheSizeGB}, para que los diez procesos quepan en la memoria de Docker Desktop.
    \item Solo \texttt{mongos} se publica hacia el equipo, en el puerto 27030, porque el 27017 lo ocupa el MongoDB instalado en Windows.
    \item En \texttt{mongos} se montan las carpetas del repositorio con los scripts (\texttt{/scripts}) y con los datos (\texttt{/datos}), para ejecutarlos con las herramientas de ese contenedor.
    \item Cada nodo de datos tiene su propio volumen, de modo que la información sobrevive si se recrea el contenedor.
\end{itemize}

\subsection{Levantamiento de los contenedores}

\begin{lstlisting}[style=plainStyle]
docker compose up -d
docker compose ps
\end{lstlisting}

\captura{01-compose.png}{Creación y estado de los once contenedores.}{1.0}

Compose creó la red, los diez volúmenes y los once contenedores. En este punto los nodos están encendidos, pero todavía no forman un cluster.

\subsection{Inicialización de los replica sets}

El replica set del config server se inicializa con el comando \texttt{replSetInitiate}. El script se ejecuta desde el contenedor de \texttt{mongos} con \texttt{--nodb} y abre la conexión al config server por su nombre en la red de Docker.

\lstinputlisting[style=jsStyle]{../configuracion/01-init-configsvr.js}

\begin{lstlisting}[style=plainStyle]
docker exec mongos mongosh --nodb --quiet --file /scripts/configuracion/01-init-configsvr.js
\end{lstlisting}

\captura{02-init-configsvr.png}{Inicialización del replica set del config server.}{1.0}

Para cada shard, \texttt{replSetInitiate} se ejecuta en su primer miembro con la lista de los tres. El primer miembro lleva prioridad 2 para que sea elegido PRIMARY; los otros dos quedan como SECONDARY y copian sus datos.

\lstinputlisting[style=jsStyle]{../configuracion/02-init-shards.js}

\begin{lstlisting}[style=plainStyle]
docker exec mongos mongosh --nodb --quiet --file /scripts/configuracion/02-init-shards.js
\end{lstlisting}

\captura{03-init-shards.png}{Inicialización de los replica sets de los tres shards.}{1.0}

Cada shard quedó con un miembro PRIMARY y dos SECONDARY. En shard3 la primera elección la ganó shard3-3; como shard3-1 tiene mayor prioridad, unos segundos después, ya sincronizado, toma el lugar de PRIMARY.

\subsection{Registro de los shards}

Se registran shard1 y shard2. El shard3 ya está inicializado, pero se agrega más adelante, con la colección ya cargada, para mostrar el balanceo de carga (sección \ref{sec:balanceo}).

\lstinputlisting[style=jsStyle]{../configuracion/03-add-shards.js}

\begin{lstlisting}[style=plainStyle]
docker exec mongos mongosh --quiet --file /scripts/configuracion/03-add-shards.js
\end{lstlisting}

\captura{04-add-shards.png}{Registro de shard1 y shard2.}{1.0}

Cada shard se registra con el nombre de su replica set y sus tres miembros, de modo que \texttt{mongos} sabe a qué nodo dirigirse si cambia el PRIMARY.

% ============================================================
\section{Utilizar la base de datos sample\_mflix y la colección movies}

\subsection{Importación de los datos}

El archivo \texttt{datos/movies.json} contiene las películas de la colección, un documento por línea en formato Extended JSON, que conserva los tipos de MongoDB como ObjectId y fechas. Se importó a través de \texttt{mongos} con \texttt{mongoimport}:

\begin{lstlisting}[style=plainStyle]
docker exec mongos mongoimport --db sample_mflix --collection movies --drop --file /datos/movies.json
\end{lstlisting}

\captura{05-importar-movies.png}{Importación de la colección movies.}{1.0}

\subsection{Monitoreo de la distribución}

Para revisar cómo se reparten los documentos se escribió un script de monitoreo que se ejecuta después de cada paso y de cada operación CRUD. Mientras la colección no está fragmentada, usa la etapa \texttt{\$collStats} para mostrar en qué shard están los documentos. Ya fragmentada, usa \texttt{\$shardedDataDistribution}, que reporta por shard los documentos que le pertenecen, su tamaño y los documentos huérfanos; además cuenta los chunks de cada shard y muestra el estado del balanceador.

\lstinputlisting[style=jsStyle]{../monitoreo/distribucion.js}

\begin{lstlisting}[style=plainStyle]
docker exec mongos mongosh --quiet --file /scripts/monitoreo/distribucion.js
\end{lstlisting}

\captura{06-distribucion-inicial.png}{Distribución de movies después de la importación.}{1.0}

Como la colección todavía no está fragmentada, las 21\,349 películas están en un solo shard, shard2, que es el shard primario de la base sample\_mflix.

% ============================================================
\section{Agregar a los documentos de películas un id\_movie}

Las películas de sample\_mflix solo tienen el \_id de tipo ObjectId. Se agregó a cada una un campo id\_movie numérico y consecutivo, que después se usa como shard key. La etapa \texttt{\$setWindowFields} numera los documentos ordenados por \_id con \texttt{\$documentNumber}, y \texttt{\$merge} guarda el número en la misma colección.

\lstinputlisting[style=jsStyle]{../configuracion/05-id-movie.js}

\begin{lstlisting}[style=plainStyle]
docker exec mongos mongosh --quiet --file /scripts/configuracion/05-id-movie.js
\end{lstlisting}

\captura{07-id-movie.png}{Asignación de id\_movie a las películas.}{1.0}

Las 21\,349 películas recibieron un id\_movie del 1 al 21\,349, sin valores repetidos.

% ============================================================
\section{Definir el Shard Key de la colección Tipo: ``Hashed''}

La shard key elegida fue \texttt{\{ id\_movie: "hashed" \}}. Los criterios fueron los siguientes:

\begin{itemize}[itemsep=2pt, topsep=4pt]
    \item Cardinalidad y frecuencia: id\_movie es único, así que tiene tantos valores como películas y ninguno se repite.
    \item Crecimiento monótono: id\_movie es consecutivo. Con sharding por rangos, las películas nuevas caerían siempre en el último chunk y en un solo shard; el hash reparte los valores consecutivos entre todos los shards.
    \item Consultas: buscar, actualizar o eliminar una película por su id\_movie se envía a un solo shard.
\end{itemize}

La colección pesa alrededor de 30 MB y el tamaño de chunk por omisión es de 128 MB. Con ese valor la colección completa cabría en un solo chunk y el balanceador no movería nada, por lo que se redujo a 1 MB. Como la colección ya tiene datos, antes de fragmentarla se crea el índice hashed sobre id\_movie.

\lstinputlisting[style=jsStyle]{../configuracion/06-shard-collection.js}

\begin{lstlisting}[style=plainStyle]
docker exec mongos mongosh --quiet --file /scripts/configuracion/06-shard-collection.js
\end{lstlisting}

\captura{08-shard-collection.png}{Definición de la shard key hashed.}{1.0}

Al fragmentar una colección que ya tiene datos, todos los documentos siguen en el shard donde se importaron, y el balanceador empieza a mover rangos hacia el otro shard. Después de unos minutos se ejecutó de nuevo el monitoreo:

\captura{09-distribucion-fragmentada.png}{Distribución de movies entre shard1 y shard2 después de fragmentar.}{1.0}

Las películas quedaron repartidas entre shard1 y shard2 en partes similares. El balanceador se detiene cuando la diferencia de datos entre los shards es menor a tres veces el tamaño de chunk, por eso el reparto no es exactamente a la mitad. La columna de huérfanos cuenta las copias que quedan en el shard de origen después de migrar un chunk: \texttt{mongos} nunca las devuelve en las consultas, y el shard las borra por sí solo unos minutos después de la migración.

% ============================================================
\section{Replica sets de los shards}
\label{sec:replicacion}

\subsection{Estado de los replica sets}

Para cada shard se consulta \texttt{replSetGetStatus} conectando directo a uno de sus miembros (no a \texttt{mongos}, que no expone el estado interno del replica set), con \texttt{directConnection} para no depender de que ese miembro sea justo el PRIMARY.

\lstinputlisting[style=jsStyle]{../replicacion/01-estado-replicasets.js}

\begin{lstlisting}[style=plainStyle]
docker exec mongos mongosh --quiet --file /scripts/replicacion/01-estado-replicasets.js
\end{lstlisting}

\captura{10-estado-replicasets.png}{Estado de los tres replica sets: un PRIMARY y dos SECONDARY cada uno.}{1.0}

Los tres shards quedaron con un PRIMARY y dos SECONDARY, con el mismo \texttt{optimeDate} entre los miembros de cada uno: no hay retraso de replicación perceptible.

\subsection{Los SECONDARY tienen los mismos datos}

Se conecta directo al PRIMARY y a un SECONDARY de shard1 y se compara el número de documentos de \texttt{sample\_mflix.movies} en cada uno. Por defecto un SECONDARY rechaza lecturas para evitar que una aplicación se conecte ahí por error, así que hay que permitirlas explícitamente con \texttt{setReadPref("secondary")}.

\lstinputlisting[style=jsStyle]{../replicacion/02-lectura-secundario.js}

\begin{lstlisting}[style=plainStyle]
docker exec mongos mongosh --quiet --file /scripts/replicacion/02-lectura-secundario.js
\end{lstlisting}

\captura{11-lectura-secundario.png}{El SECONDARY de shard1 tiene copiados los mismos documentos que su PRIMARY.}{0.9}

El SECONDARY devolvió el mismo total que el PRIMARY, y el documento de ejemplo salió idéntico en ambos: la replicación dentro del shard está al día.

\subsection{Prueba de failover}

Se detiene el contenedor del PRIMARY de shard1 (\texttt{shard1-1}) directamente desde Docker, sin tocar mongosh, y se espera unos segundos a que los dos SECONDARY elijan uno nuevo.

\begin{lstlisting}[style=plainStyle]
docker stop shard1-1
\end{lstlisting}

\lstinputlisting[style=jsStyle]{../replicacion/03-failover-caida.js}

\begin{lstlisting}[style=plainStyle]
docker exec mongos mongosh --quiet --file /scripts/replicacion/03-failover-caida.js
\end{lstlisting}

\captura{12-failover-caida.png}{Nuevo PRIMARY elegido en shard1 y consulta a través de mongos funcionando sin shard1-1.}{0.9}

Con \texttt{shard1-1} caído, \texttt{shard1-2} fue elegido PRIMARY y el cluster, visto desde \texttt{mongos}, siguió respondiendo con el total completo de películas: el router no necesita saber cuál nodo específico está activo, solo que el shard tenga un PRIMARY.

Después se vuelve a encender \texttt{shard1-1}. Como tiene prioridad más alta, en cuanto termina de ponerse al día provoca una nueva elección y recupera el PRIMARY.

\begin{lstlisting}[style=plainStyle]
docker start shard1-1
\end{lstlisting}

\lstinputlisting[style=jsStyle]{../replicacion/04-recuperacion.js}

\begin{lstlisting}[style=plainStyle]
docker exec mongos mongosh --quiet --file /scripts/replicacion/04-recuperacion.js
\end{lstlisting}

\captura{13-failover-recuperacion.png}{shard1-1 vuelve como SECONDARY y, al ponerse al día, recupera el PRIMARY por su prioridad.}{0.9}

\texttt{shard1-1} reapareció primero como SECONDARY (tenía que replicar lo que pasó mientras estuvo apagado) y segundos después, ya sincronizado, volvió a ser PRIMARY — exactamente el comportamiento esperado de una prioridad más alta.

% ============================================================
\section{Mostrar un ejemplo de balanceo de carga entre shards}
\label{sec:balanceo}

Antes de agregar shard3, la colección ya fragmentada entre shard1 y shard2 se revisa con el script de monitoreo.

\begin{lstlisting}[style=plainStyle]
docker exec mongos mongosh --quiet --file /scripts/monitoreo/distribucion.js
\end{lstlisting}

\captura{15-antes-balanceo.png}{Distribución entre shard1 y shard2 antes de registrar shard3.}{1.0}

\texttt{shard3} ya estaba levantado e inicializado como replica set desde el principio (\texttt{configuracion/02-init-shards.js}); lo único que faltaba era registrarlo en el cluster con \texttt{sh.addShard}.

\lstinputlisting[style=jsStyle]{../balanceo/01-agregar-shard3.js}

\begin{lstlisting}[style=plainStyle]
docker exec mongos mongosh --quiet --file /scripts/balanceo/01-agregar-shard3.js
\end{lstlisting}

\captura{16-agregar-shard3.png}{shard3 registrado; el cluster ya tiene tres shards.}{0.9}

En cuanto \texttt{mongos} ve disponible al shard nuevo, el balanceador empieza a moverle chunks por su cuenta, sin que haya que pedírselo. El avance se sigue corriendo el mismo script de monitoreo varias veces.

\begin{lstlisting}[style=plainStyle]
docker exec mongos mongosh --quiet --file /scripts/monitoreo/distribucion.js
\end{lstlisting}

\captura{17-durante-balanceo.png}{Migración en curso: shard3 ya tiene documentos y los otros dos van bajando.}{1.0}

Unos minutos después, con la ronda de balanceo terminada, los tres shards quedan con una parte similar de la colección.

\captura{18-despues-balanceo.png}{Distribución final entre los tres shards, alrededor de un tercio cada uno.}{1.0}

\begin{lstlisting}[style=plainStyle]
docker exec mongos mongosh --quiet --eval "sh.status()"
\end{lstlisting}

\captura{19-sh-status.png}{Estado final del cluster con los tres shards registrados.}{1.0}

El reparto no queda exactamente a un tercio cada uno porque el balanceador deja de mover datos en cuanto la diferencia entre shards baja de tres veces el tamaño de chunk (3\,MB); aun así, los tres terminan con una proporción parecida de la colección.

% ============================================================
\section{Realizar operaciones CRUD sobre la base de datos distribuida}

% ------------------------------------------------------------
% SECCIÓN A CARGO DE OSVALDO. Ver CONTEXTO.md.
% Scripts en crud/, capturas 20-28.
% ------------------------------------------------------------

Las operaciones CRUD se realizaron después del balanceo, con los tres shards registrados, y siempre a través de \texttt{mongos}. Las películas que se crean en esta sección tienen un id\_movie mayor a 21\,349 para distinguirlas de las originales, que no se modifican. Al terminar, la colección vuelve a tener 21\,349 películas.

\subsection{Funciones de apoyo y monitoreo}

Los cuatro scripts cargan el archivo \texttt{crud/comun.js}, que define dos funciones de monitoreo. La función \texttt{distribucion()} usa \texttt{\$shardedDataDistribution} para mostrar cuántos documentos tiene cada shard justo después de cada operación individual. La función \texttt{destino()} usa \texttt{explain()} para mostrar el plan de \texttt{mongos}: \texttt{SINGLE\_SHARD} cuando la operación va a un solo shard y \texttt{SHARD\_MERGE} cuando consulta a todos y junta los resultados, junto con los documentos que devolvió cada shard.

\lstinputlisting[style=jsStyle]{../crud/comun.js}

Además, después de cada script se ejecutó el script de monitoreo completo, \texttt{distribucion.js}. Antes de comenzar, la distribución era la que dejó el balanceo:

\begin{lstlisting}[style=plainStyle]
docker exec mongos mongosh --quiet --file /scripts/monitoreo/distribucion.js
\end{lstlisting}

\captura{20-distribucion-antes-crud.png}{Distribución de movies antes de las operaciones CRUD.}{1.0}

\subsection{Create}

Se insertó una película con \texttt{insertOne} (id\_movie 21\,350) y un lote de 1\,000 películas con \texttt{insertMany} (id\_movie 21\,351 a 22\,350). Todas incluyen id\_movie, porque es la shard key: \texttt{mongos} calcula el hash de su valor para decidir en qué shard se guarda cada documento.

\lstinputlisting[style=jsStyle]{../crud/create.js}

\begin{lstlisting}[style=plainStyle]
docker exec mongos mongosh sample_mflix --quiet --file /scripts/crud/create.js
\end{lstlisting}

\captura{21-create.png}{Inserción de películas con insertOne e insertMany.}{1.0}

La película de \texttt{insertOne} se guardó en un solo shard, el que corresponde al hash de su id\_movie, y solo ese shard aumentó en un documento. Aunque los id\_movie del lote son consecutivos, el hash repartió las 1\,000 películas entre los tres shards, en proporción a la parte del rango de hash que le pertenece a cada uno.

\begin{lstlisting}[style=plainStyle]
docker exec mongos mongosh --quiet --file /scripts/monitoreo/distribucion.js
\end{lstlisting}

\captura{22-monitoreo-create.png}{Distribución de movies después de las inserciones.}{1.0}

El total de la colección aumentó a 22\,350 documentos. Además, entre la inserción y este monitoreo el balanceador movió por su cuenta un chunk de 661 documentos de shard1 a shard2: shard1 bajó de 8\,791 a 8\,130 documentos, shard2 pasó de 10 a 11 chunks y los huérfanos de shard1 aumentaron justo en 661. Antes de las inserciones, shard1 tenía casi 3\,MB más que los otros dos shards; las películas nuevas hicieron que la diferencia superara el umbral de tres veces el tamaño de chunk y el balanceador volvió a repartir los datos sin intervención.

\subsection{Read}

Se realizaron tres consultas: por id\_movie, por campos que no forman parte de la shard key (genres, year e imdb.rating) y por un rango de id\_movie que corresponde al lote insertado.

\lstinputlisting[style=jsStyle]{../crud/read.js}

\begin{lstlisting}[style=plainStyle]
docker exec mongos mongosh sample_mflix --quiet --file /scripts/crud/read.js
\end{lstlisting}

\captura{23-read.png}{Consultas sobre la colección distribuida con explain().}{1.0}

Las búsquedas por id\_movie usan el plan \texttt{SINGLE\_SHARD}: \texttt{mongos} calcula el hash del valor y consulta únicamente el shard que tiene ese documento. La consulta por genres, year e imdb.rating no incluye la shard key, así que \texttt{mongos} la envía a los tres shards y junta sus resultados (\texttt{SHARD\_MERGE}). Lo mismo ocurre con el rango de id\_movie: con una shard key hashed los valores consecutivos están repartidos, por lo que un rango también se consulta en todos los shards.

\begin{lstlisting}[style=plainStyle]
docker exec mongos mongosh --quiet --file /scripts/monitoreo/distribucion.js
\end{lstlisting}

\captura{24-monitoreo-read.png}{Distribución de movies después de las consultas.}{1.0}

Las consultas no modifican documentos, por lo que la distribución quedó igual que después de las inserciones.

\subsection{Update}

Se realizaron tres actualizaciones: un \texttt{updateOne} por id\_movie que modifica la calificación y los géneros de la película insertada, un \texttt{updateMany} que marca como revisadas las 1\,001 películas insertadas, y un \texttt{updateOne} que cambia el valor de id\_movie de una película.

\lstinputlisting[style=jsStyle]{../crud/update.js}

\begin{lstlisting}[style=plainStyle]
docker exec mongos mongosh sample_mflix --quiet --file /scripts/crud/update.js
\end{lstlisting}

\captura{25-update.png}{Actualización de películas con updateOne y updateMany.}{1.0}

El \texttt{updateOne} por id\_movie se envió a un solo shard y el \texttt{updateMany} a los tres; en ambos casos \texttt{matchedCount} y \texttt{modifiedCount} coinciden con los documentos esperados. Cambiar campos que no son la shard key no mueve documentos, por lo que la distribución no cambió. En cambio, al cambiar el valor de id\_movie, el nuevo hash pertenece a otro shard y MongoDB movió el documento: el shard de origen perdió una película y el de destino ganó una, sin cambiar el total.

\begin{lstlisting}[style=plainStyle]
docker exec mongos mongosh --quiet --file /scripts/monitoreo/distribucion.js
\end{lstlisting}

\captura{26-monitoreo-update.png}{Distribución de movies después de las actualizaciones.}{1.0}

\subsection{Delete}

Por último, se eliminó la película de \texttt{insertOne} con \texttt{deleteOne}, usando su id\_movie actual, y el resto del lote con \texttt{deleteMany}.

\lstinputlisting[style=jsStyle]{../crud/delete.js}

\begin{lstlisting}[style=plainStyle]
docker exec mongos mongosh sample_mflix --quiet --file /scripts/crud/delete.js
\end{lstlisting}

\captura{27-delete.png}{Eliminación de películas con deleteOne y deleteMany.}{1.0}

El \texttt{deleteOne} se envió solo al shard que tenía la película, mientras que el \texttt{deleteMany} se ejecutó en los tres shards y cada uno borró sus propios documentos del lote, 1\,000 en total.

\begin{lstlisting}[style=plainStyle]
docker exec mongos mongosh --quiet --file /scripts/monitoreo/distribucion.js
\end{lstlisting}

\captura{28-monitoreo-delete.png}{Distribución de movies después de las eliminaciones.}{1.0}

La colección volvió a tener 21\,349 películas. El reparto final no es igual al de la figura de la distribución antes de las operaciones CRUD: shard1 tiene 631 películas menos y shard2 631 más. Son las películas originales que venían en el chunk que el balanceador movió después de las inserciones (los otros 30 documentos de ese chunk eran películas nuevas y se eliminaron con el \texttt{deleteMany}). Las películas siguen completas; solo cambió el shard que guarda una parte de ellas.

% ============================================================
\section{Conclusión}

% Al final, entre todos: dos o tres párrafos con lo que se observó
% (replicación, balanceo, distribución en las operaciones CRUD).
Con este proyecto se implementó una base de datos distribuida con MongoDB sobre Docker: un config server, tres shards configurados como replica sets de tres miembros y un router \texttt{mongos}, con la colección movies de sample\_mflix fragmentada por una shard key hashed sobre id\_movie. Para las aplicaciones, el cluster se comporta como una sola base de datos: todas las operaciones se hicieron a través de \texttt{mongos}, que decide en qué shard está cada documento.

La replicación dio tolerancia a fallos dentro de cada shard. Los SECONDARY mantuvieron los mismos datos que su PRIMARY, y al detener el PRIMARY de shard1 se eligió otro miembro automáticamente sin que las consultas a través de \texttt{mongos} dejaran de funcionar. Al encenderlo de nuevo, el nodo se puso al día y recuperó su papel por su prioridad.

El sharding repartió la carga entre los shards. Al registrar shard3 con la colección ya cargada, el balanceador le movió chunks hasta dejar a los tres con una parte similar de los datos, y volvió a actuar por sí solo durante las operaciones CRUD, cuando las inserciones desbalancearon a shard1. La shard key hashed repartió entre los tres shards incluso valores consecutivos de id\_movie, y el monitoreo con \texttt{explain()} mostró su efecto en las consultas: las operaciones que filtran por id\_movie se envían a un solo shard, mientras que las que usan otros campos, o rangos de id\_movie, se ejecutan en todos. Por eso la elección de la shard key debe considerar tanto el reparto de los datos como las consultas más frecuentes de la aplicación.

% ============================================================
\section*{Referencias}
\addcontentsline{toc}{section}{Referencias}

\begin{list}{}{%
    \setlength{\leftmargin}{1.27cm}
    \setlength{\itemindent}{-1.27cm}
    \setlength{\itemsep}{6pt}
    \setlength{\parsep}{0pt}}
\item Borucki, A. (2025). Using replication and sharding. En \textit{MongoDB 8.0 in Action} (3a ed., cap. 9). Manning.
\item MongoDB, Inc. (s.f.). \textit{Replication}. MongoDB Manual. \url{https://www.mongodb.com/docs/manual/replication/}
\item MongoDB, Inc. (s.f.). \textit{Sharding}. MongoDB Manual. \url{https://www.mongodb.com/docs/manual/sharding/}
\item udaykirankavaturu. (s.f.). \textit{sharding-demo} [Repositorio]. GitHub. \url{https://github.com/udaykirankavaturu/sharding-demo}
\end{list}

\end{document}
